% ============================================================
% APPENDIX A — PROOFS
% ============================================================
\section{Proofs}
\label{app:proofs}

We collect the proofs of all stated results.  For clarity, we
follow the notation of \citet{xu2023conformal} and explicitly
identify where the seasonal structure introduces differences.
Without loss of generality, we prove results for $t = T+1$
(one-step-ahead prediction); extensions to $t > T+1$ hold by
the same argument whenever Assumptions~\ref{ass:strat_a1}--\ref{ass:strat_a2}
are satisfied at indices $t-T, \ldots, t$.

%-----------------------------------------------------------------
\subsection{Proof of Lemma~\ref{lem:bias_properties}}
\label{app:proofs:lemma_bias}
%-----------------------------------------------------------------

Part~(i): If $\sigma_s = \sigma$ for all $s$, then
$F_s(x) = G(x/\sigma)$ for every $s$, so
$F_{\sstar}(x) = \bar{F}(x)$ pointwise and
$\bsstar = \sup_x|F_{\sstar}(x) - \bar{F}(x)| = 0$.

Part~(ii): By Definition~\ref{def:seasonal_bias},
$\bsstar = \sup_x|G(x/\sigma_{\sstar}) - (1/\Sper)\sum_s G(x/\sigma_s)|$.
Neither $G$, $\sigma_1,\ldots,\sigma_\Sper$, nor $\Sper$ depend on
the sample size $T$.  Hence $\bsstar$ does not depend on $T$.
\qed

%-----------------------------------------------------------------
\subsection{An Auxiliary Concentration Inequality}
\label{app:proofs:dkw_noniid}
%-----------------------------------------------------------------

The sharp DKW inequality of \citet{massart1990tight} applies to
i.i.d.\ samples.  The pooled empirical CDF under
Assumption~\ref{ass:seasonal} aggregates independent but
\emph{not} identically distributed errors, so we record the
following standard substitute, proved by symmetrisation and the
bounded-differences inequality
\citep{boucheron2013concentration}.

\begin{lemma}[Uniform concentration, independent non-i.i.d.\ case]
\label{lem:dkw_noniid}
Let $\epsilon_1, \ldots, \epsilon_T$ be independent random
variables with (possibly distinct) CDFs $F_{s(1)}, \ldots,
F_{s(T)}$, let
$\tilde{F}_T(x) = T^{-1}\sum_{t=1}^T \Ind{\epsilon_t \leq x}$
and $\bar{F}_T(x) = T^{-1}\sum_{t=1}^T F_{s(t)}(x)$.  Then
\[
  \E\Bigl[\sup_x \bigl|\tilde{F}_T(x) - \bar{F}_T(x)\bigr|\Bigr]
  \;\leq\; 2\sqrt{\frac{2\log(2(T+1))}{T}},
  \qquad
  \Prob\Bigl(\sup_x \bigl|\tilde{F}_T(x) - \bar{F}_T(x)\bigr|
  > \E[\,\cdot\,] + u\Bigr) \;\leq\; e^{-2Tu^2}
\]
for every $u > 0$.  In particular, with probability at least
$1 - (16T)^{-1}$,
\[
  \sup_x \bigl|\tilde{F}_T(x) - \bar{F}_T(x)\bigr|
  \;\leq\;
  4\sqrt{\frac{\log(16T)}{T}}.
\]
\end{lemma}

\begin{proof}
Replacing $\epsilon_t$ by an independent copy changes
$\sup_x|\tilde{F}_T(x) - \bar{F}_T(x)|$ by at most $1/T$, so the
bounded-differences inequality gives the stated tail bound around
the mean.  For the mean, symmetrisation yields
$\E[\sup_x|\tilde{F}_T - \bar{F}_T|]
\leq 2\,\E\bigl[\sup_x |T^{-1}\sum_t \sigma_t
\Ind{\epsilon_t \leq x}|\bigr]$
for i.i.d.\ Rademacher signs $\{\sigma_t\}$; for any fixed sample,
as $x$ varies the vector
$(\Ind{\epsilon_1 \leq x}, \ldots, \Ind{\epsilon_T \leq x})$
takes at most $T+1$ distinct values, each with Euclidean norm at
most $\sqrt{T}$, so the maximal inequality for finite classes
\citep[see][]{boucheron2013concentration}
bounds the Rademacher average by $\sqrt{2\log(2(T+1))/T}$.
For the final display, take $u = \sqrt{\log(16T)/(2T)}$, so that
$e^{-2Tu^2} = (16T)^{-1}$, and note that
$\log(2(T+1)) \leq \log(16T)$ for all $T \geq 1$; hence
$\E[\sup] + u \leq
(2\sqrt{2} + 1/\sqrt{2})\sqrt{\log(16T)/T}
\leq 4\sqrt{\log(16T)/T}$.
\end{proof}

%-----------------------------------------------------------------
\subsection{Proof of Proposition~\ref{prop:pooled_gap}}
\label{app:proofs:prop_pooled}
%-----------------------------------------------------------------

We follow the proof of \citet[Theorem~1]{xu2023conformal}, modified
to account for the heterogeneous seasonal distributions.

\medskip
\noindent\textit{Step 1 (p-value equivalence).}
Define the stratified empirical $p$-value using the pooled residuals:
\[
  \hat{p}_{T+1}^{\text{pool}}
  \;:=\;
  \hat{F}_T^{\text{pool}}(\hat{\epsilon}_{T+1}),
  \qquad
  \hat{F}_T^{\text{pool}}(x)
  = \frac{1}{T}\sum_{t=1}^T \Ind{\hat{\epsilon}_t \leq x}.
\]
The coverage event satisfies:
\begin{equation}
  \label{eq:pval_equiv}
  Y_{T+1} \in \hat{C}_{T+1}^{\alpha,\text{pool}}
  \;\iff\;
  \hat{\beta} \leq \hat{p}_{T+1}^{\text{pool}} \leq 1-\alpha+\hat{\beta}
\end{equation}
(cf.\ \citealt{xu2023conformal}, Sec.~3.1).  Thus:
\begin{align}
  \label{eq:gap_start}
  &\bigl|\Prob(Y_{T+1} \in \hat{C}_{T+1}^{\alpha,\text{pool}} \mid X_{T+1}=x_{T+1}) - (1-\alpha)\bigr|
  \notag\\
  &= \bigl|\Prob(\hat{\beta} \leq \hat{F}_T^{\text{pool}}(\hat{\epsilon}_{T+1}) \leq 1-\alpha+\hat{\beta})
     - \Prob(\beta \leq F_{\sstar}(\epsilon_{T+1}) \leq 1-\alpha+\beta)\bigr|,
\end{align}
where we used the Probability Integral Transform:
$F_{\sstar}(\epsilon_{T+1}) \sim \mathrm{Unif}[0,1]$ since
$\epsilon_{T+1} \sim F_{\sstar}$ (continuous), and thus
$\Prob(\beta \leq F_{\sstar}(\epsilon_{T+1}) \leq 1-\alpha+\beta)
= 1-\alpha$ for any $\beta \in [0,\alpha]$.

\medskip
\noindent\textit{Step 2 (three-step triangle decomposition).}
Introduce two auxiliary CDFs:
\begin{align}
  \tilde{F}_T^{\text{pool}}(x)
  &= \frac{1}{T}\sum_{t=1}^T \Ind{\epsilon_t \leq x}
  \quad\text{(pooled empirical CDF of true errors)}, \notag\\
  \bar{F}(x)
  &= \frac{1}{\Sper}\sum_{s=1}^\Sper F_s(x)
  \quad\text{(pooled mixture CDF, defined in~\eqref{eq:pooled_cdf})}. \notag
\end{align}
By the triangle inequality:
\begin{equation}
  \label{eq:three_step}
  \bigl|\hat{F}_T^{\text{pool}}(x) - F_{\sstar}(x)\bigr|
  \;\leq\;
  \underbrace{\bigl|\hat{F}_T^{\text{pool}}(x) - \tilde{F}_T^{\text{pool}}(x)\bigr|}_{(A)}
  +
  \underbrace{\bigl|\tilde{F}_T^{\text{pool}}(x) - \bar{F}(x)\bigr|}_{(B)}
  +
  \underbrace{\bigl|\bar{F}(x) - F_{\sstar}(x)\bigr|}_{(C)}.
\end{equation}

\medskip
\noindent\textit{Step 3 (bounding term (A) via LOO coupling).}
By Lemma~\ref{lem:coupling}, $\hat{\epsilon}_t = \epsilon_t +
(f(x_t) - \hat{f}_{-t}^{\,\phi}(x_t))$.  We apply Lemma~2 of
\citet{xu2023conformal}, whose proof extends to non-identically
distributed errors once the common CDF $F$ in step~(iii) of their
argument is replaced by the mixture $\bar{F}$, which is Lipschitz
with constant $L = \max_s L_s$ (Assumption~\ref{ass:xu_a2} is used
unchanged):
\[
  \sup_x\bigl|\hat{F}_T^{\text{pool}}(x) - \tilde{F}_T^{\text{pool}}(x)\bigr|
  \;\leq\;
  (L+1)\deltaT^{2/3} + 2\sup_x\bigl|\tilde{F}_T^{\text{pool}}(x) - \bar{F}(x)\bigr|,
\]
where $L = \max_s L_s$.

\medskip
\noindent\textit{Step 4 (bounding term (B) via
Lemma~\ref{lem:dkw_noniid}).}
The errors $\{\epsilon_t\}_{t=1}^T$ are mutually independent
with marginal distributions $F_{s(t)}$
(Assumption~\ref{ass:seasonal}), and under the balanced panel
$E[\tilde{F}_T^{\text{pool}}(x)]
= T^{-1}\sum_{t=1}^T F_{s(t)}(x) = \bar{F}(x)$.
Since the errors are not identically distributed, the sharp DKW
inequality of \citet{massart1990tight} is not available; we
invoke Lemma~\ref{lem:dkw_noniid} instead.  Define the good event
\[
  A_T \;=\;
  \Bigl\{\sup_x\bigl|\tilde{F}_T^{\text{pool}}(x)-\bar{F}(x)\bigr|
  \leq 4\sqrt{\log(16T)/T}\Bigr\},
  \qquad
  \Prob(A_T^c) \;\leq\; (16T)^{-1} \;\leq\; \sqrt{\log(16T)/T}.
\]
On the event $A_T$, term~(B) is bounded by
$4\sqrt{\log(16T)/T}$.

\medskip
\noindent\textit{Step 5 (term (C)).}
By Definition~\ref{def:seasonal_bias}:
$\sup_x|\bar{F}(x) - F_{\sstar}(x)| = \bsstar$.
This term is independent of both $T$ and the realisation of the data.

\medskip
\noindent\textit{Step 6 (assembling the bound and identifying the bias term).}

Ignoring estimation error for the moment (it contributes the $(L+1)(\deltaT^{2/3}+\deltaT)$ terms via the argument in the proof of Theorem~1 of \citealt{xu2023conformal}, pp.~31--33 of the arXiv version), the essential coverage computation is:
\begin{align}
  &\Prob\bigl(Y_{T+1} \in \hat{C}_{T+1}^{\alpha,\text{pool}}\bigr)
  \;\approx\;
  \Prob\bigl(Q_{\hat\beta}(\hat{\boldsymbol\epsilon}^{\text{pool}})
             \leq \epsilon_{T+1} \leq
             Q_{1-\alpha+\hat\beta}(\hat{\boldsymbol\epsilon}^{\text{pool}})\bigr)
  \notag\\
  &\approx\;
  F_{\sstar}\!\bigl(\bar{F}^{-1}(1-\alpha+\hat\beta)\bigr)
  - F_{\sstar}\!\bigl(\bar{F}^{-1}(\hat\beta)\bigr),
  \label{eq:cov_approx}
\end{align}
where the second approximation uses the fact that
$Q_p(\hat{\boldsymbol\epsilon}^{\text{pool}}) \approx \bar{F}^{-1}(p)$ as
the pooled empirical CDF concentrates around $\bar{F}$.
The coverage gap relative to $(1-\alpha)$ is:
\begin{align*}
  &F_{\sstar}\!\bigl(\bar{F}^{-1}(1-\alpha+\hat\beta)\bigr)
   - F_{\sstar}\!\bigl(\bar{F}^{-1}(\hat\beta)\bigr)
   - (1-\alpha) \\
  &=\bigl[F_{\sstar}\!\bigl(\bar{F}^{-1}(1-\alpha+\hat\beta)\bigr)
      - \bar{F}\!\bigl(\bar{F}^{-1}(1-\alpha+\hat\beta)\bigr)\bigr]
   - \bigl[F_{\sstar}\!\bigl(\bar{F}^{-1}(\hat\beta)\bigr)
      - \bar{F}\!\bigl(\bar{F}^{-1}(\hat\beta)\bigr)\bigr].
\end{align*}
Since $\bar{F}(\bar{F}^{-1}(u)) = u$ for any $u \in (0,1)$, each
bracket has absolute value at most
$\sup_x|F_{\sstar}(x) - \bar{F}(x)| = \bsstar$.
By the triangle inequality, the absolute value of the entire
expression is at most $\bsstar + \bsstar = 2\bsstar$.

Incorporating the sampling and estimation error terms
from Steps~3--5 (which follow \citealt[proof of Theorem~1]{xu2023conformal},
with the concentration input supplied by
Lemma~\ref{lem:dkw_noniid} in place of the i.i.d.\ DKW
inequality, and Assumption~\ref{ass:xu_a2}) yields:
\[
  \bigl|\Prob(Y_{T+1} \in \hat{C}_{T+1}^{\alpha,\text{pool}} \mid X_{T+1}=x_{T+1}) - (1-\alpha)\bigr|
  \;\leq\;
  c_0\sqrt{\frac{\log(16T)}{T}}
  + 4(L+1)(\deltaT^{2/3}+\deltaT)
  + 2\bsstar,
\]
where $c_0$ is the absolute constant obtained by tracking the
bound of Lemma~\ref{lem:dkw_noniid} through
\citeauthor{xu2023conformal}'s argument ($c_0 = 12$ in the
i.i.d.\ case, where the sharp DKW constant applies).  This
is~\eqref{eq:pooled_gap}. The factor $2\bsstar$ is
\emph{exact} in the sense that it arises from a difference of two
CDF evaluations, each bounded by $\bsstar$.\qed

%-----------------------------------------------------------------
\subsection{Proof of Theorem~\ref{thm:strat_coverage}}
\label{app:proofs:thm1s}
%-----------------------------------------------------------------

The proof applies the argument of
\citet[Theorem~1]{xu2023conformal} to the stratified
sub-sequence of length $\Ts$.

\medskip
\noindent\textit{Step 1 (stratified p-value equivalence).}
Define the within-stratum empirical CDF of LOO residuals:
\[
  \hat{F}_{\Ts}^{(\sstar)}(x)
  = \frac{1}{\Ts}\sum_{t \in \calT_{\sstar}} \Ind{\hat{\epsilon}_t \leq x}.
\]
The coverage event satisfies:
\[
  Y_{T+1} \in \hat{C}_{T+1}^{\alpha,\text{strat}}
  \;\iff\;
  \hat{\beta}^{(\sstar)} \leq
  \hat{F}_{\Ts}^{(\sstar)}(\hat{\epsilon}_{T+1})
  \leq 1-\alpha+\hat{\beta}^{(\sstar)},
\]
which mirrors~\eqref{eq:pval_equiv} with the stratified buffer.
Thus:
\begin{align}
  \label{eq:strat_gap_start}
  &\bigl|\Prob(Y_{T+1} \in \hat{C}_{T+1}^{\alpha,\text{strat}} \mid X_{T+1}=x_{T+1}) - (1-\alpha)\bigr|
  \notag\\
  &= \bigl|\Prob(\hat{\beta}^{(\sstar)} \leq
    \hat{F}_{\Ts}^{(\sstar)}(\hat{\epsilon}_{T+1})
    \leq 1-\alpha+\hat{\beta}^{(\sstar)})
     - \Prob(\beta \leq F_{\sstar}(\epsilon_{T+1}) \leq 1-\alpha+\beta)\bigr|.
\end{align}
The second probability equals $1-\alpha$ by the Probability
Integral Transform applied to $\epsilon_{T+1} \sim F_{\sstar}$
(Assumption~\ref{ass:strat_a1}).

\medskip
\noindent\textit{Step 2 (two-step triangle decomposition).}
Introduce the within-stratum empirical CDF of true errors:
$\tilde{F}_{\Ts}^{(\sstar)}(x)
= \Ts^{-1}\sum_{t \in \calT_{\sstar}} \Ind{\epsilon_t \leq x}$.
By the triangle inequality:
\[
  \bigl|\hat{F}_{\Ts}^{(\sstar)}(x) - F_{\sstar}(x)\bigr|
  \;\leq\;
  \underbrace{\bigl|\hat{F}_{\Ts}^{(\sstar)}(x) - \tilde{F}_{\Ts}^{(\sstar)}(x)\bigr|}_{(A')}
  +
  \underbrace{\bigl|\tilde{F}_{\Ts}^{(\sstar)}(x) - F_{\sstar}(x)\bigr|}_{(B')}.
\]
Note that the third term $(C)$ from the pooled proof is absent:
the stratified empirical CDF compares directly with $F_{\sstar}$,
not with $\bar{F}$.

\medskip
\noindent\textit{Step 3 (bounding term $(A')$ via LOO coupling).}
The within-stratum LOO residuals $\hat{\epsilon}_t$ for
$t \in \calT_{\sstar}$ satisfy the coupling~\eqref{eq:coupling}
with the same pointwise estimation error.  Applying Lemma~2 of
\citet{xu2023conformal} with the within-stratum quantities
from Assumption~\ref{ass:strat_a2}:
\[
  \sup_x\bigl|\hat{F}_{\Ts}^{(\sstar)}(x) - \tilde{F}_{\Ts}^{(\sstar)}(x)\bigr|
  \;\leq\;
  (L_{\sstar}+1)\deltaTs^{2/3}
  + 2\sup_x\bigl|\tilde{F}_{\Ts}^{(\sstar)}(x) - F_{\sstar}(x)\bigr|.
\]

\medskip
\noindent\textit{Step 4 (bounding term $(B')$ via DKW).}
The $\Ts$ within-stratum errors $\{\epsilon_t : t \in \calT_{\sstar}\}$
are i.i.d.\ from $F_{\sstar}$ by Assumption~\ref{ass:strat_a1}.
Applying the DKW inequality \citep{massart1990tight} directly:
$\Prob(\sup_x|\tilde{F}_{\Ts}^{(\sstar)}(x)-F_{\sstar}(x)| > u)
\leq 2e^{-2\Ts u^2}$.
Following the proof of Lemma~1 in \citet{xu2023conformal}, choose
$u = \sqrt{W(16\Ts)}/(2\sqrt{\Ts}) \leq \sqrt{\log(16\Ts)/\Ts}$,
where $W(\cdot)$ denotes the Lambert function, defined by
$W(z)e^{W(z)} = z$.
Define the event $A_T^{(\sstar)}$ on which
$\sup_x|\tilde{F}_{\Ts}^{(\sstar)}(x)-F_{\sstar}(x)|
\leq \sqrt{\log(16\Ts)/\Ts}$;
this event satisfies
$\Prob((A_T^{(\sstar)})^c) \leq \sqrt{\log(16\Ts)/\Ts}$.

\medskip
\noindent\textit{Step 5 (assembling the stratified bound).}
Following the argument in the proof of Theorem~1 of
\citet{xu2023conformal} (pp.~31--33 of the arXiv version),
conditioning on $A_T^{(\sstar)}$ and using
$F_{\sstar}(\epsilon_{T+1}) \sim \mathrm{Unif}[0,1]$:
\[
  \Prob\!\bigl(|F_{\sstar}(\epsilon_{T+1}) - \gamma|
               \leq |\hat{F}_{\Ts}^{(\sstar)}(\hat{\epsilon}_{T+1})
                    - F_{\sstar}(\epsilon_{T+1})|
        \,\big|\, A_T^{(\sstar)}\bigr)
  \;\leq\;
  6\sqrt{\frac{\log(16\Ts)}{\Ts}}
  + 2(L_{\sstar}+1)(\deltaTs^{2/3}+\deltaTs),
\]
for $\gamma \in \{\beta, 1-\alpha+\beta\}$.  Accounting for
$\Prob\bigl((A_T^{(\sstar)})^c\bigr) \leq \sqrt{\log(16\Ts)/\Ts}$
and summing over the two values of $\gamma$, exactly as in the
original argument, together
with~\eqref{eq:strat_gap_start}, yields~\eqref{eq:strat_bound}.\qed

%-----------------------------------------------------------------
\subsection{Auxiliary: Mixing of Sub-sampled Processes}
\label{app:proofs:submixing}
%-----------------------------------------------------------------

\begin{proposition}[Strong mixing of season-$s$ sub-sequence]
\label{prop:submixing}
Let $\{\epsilon_t\}_{t \geq 1}$ be a strongly
mixing sequence with mixing coefficients $\{\alpha_n\}_{n \geq 0}$
(stationarity is not required for the mixing statement).
For any $s \in \{1,\ldots,\Sper\}$, the sub-sequence
$\{\epsilon_t : s(t) = s\}$ (observations at time indices
$s, s+\Sper, s+2\Sper, \ldots$) is strongly
mixing with mixing coefficients bounded above by
$\alpha_\Sper, \alpha_{2\Sper}, \ldots$
If, in addition, the process is strictly stationary or, more
generally, $\Sper$-periodically stationary (its joint law is
invariant under shifts by $\Sper$, as under
Assumption~\ref{ass:seasonal} with i.i.d.\ $\eta_t$), the
sub-sequence is strictly stationary.
\end{proposition}

\begin{proof}
Let $Z_k = \epsilon_{s + (k-1)\Sper}$ for $k = 1, 2, \ldots$ be
the sub-sequence.  The sub-$\sigma$-algebra generated by
$Z_1, \ldots, Z_k$ satisfies
$\sigma(Z_1, \ldots, Z_k) = \sigma(\epsilon_s, \ldots, \epsilon_{s+(k-1)\Sper})
\subseteq \sigma(\epsilon_1, \ldots, \epsilon_{s+(k-1)\Sper})$.
Similarly, $\sigma(Z_{k+n}, Z_{k+n+1}, \ldots) \subseteq
\sigma(\epsilon_{s+(k+n-1)\Sper}, \epsilon_{s+(k+n)\Sper}, \ldots)$.
The mixing coefficient of $\{Z_k\}$ at lag $n$ is therefore
bounded by the mixing coefficient of $\{\epsilon_t\}$ at lag
$(k+n)\Sper - k\Sper = n\Sper$, i.e., by $\alpha_{n\Sper}$.
Since $\alpha_{n\Sper} \to 0$ as $n \to \infty$ (because
$\alpha_n \to 0$), the sub-sequence
is strongly mixing.  If $\{\epsilon_t\}$ is strictly stationary
or $\Sper$-periodically stationary, the joint law of
$(Z_{k_1+j}, \ldots, Z_{k_m+j})$ is invariant in $j$, since a
unit shift of $\{Z_k\}$ corresponds to a shift by $\Sper$ of
$\{\epsilon_t\}$; hence $\{Z_k\}$ is strictly stationary.
\end{proof}

%-----------------------------------------------------------------
\subsection{Proof of Theorem~\ref{thm:width_bound}}
\label{app:proofs:thm2s}
%-----------------------------------------------------------------

The proof mirrors \citet[Theorem~3]{xu2023conformal},
with all global quantities ($T$, $\deltaT$, $L$, etc.) replaced
by their within-stratum counterparts
($\Ts$, $\deltaTs$, $L_{\sstar}$, etc.).

\medskip
Define $\hat{F}_{\Ts}^{(\sstar),-1}$ as the interpolated inverse
CDF of the within-stratum residuals, with Lipschitz constant
$K'_{\sstar}$ as defined in Theorem~\ref{thm:width_bound}.
The symmetric difference $\Delta_{\sstar}(T)$ between
$\hat{C}_{T+1}^{\alpha,\text{strat}}$ and the oracle $C_{T+1}^\alpha$
is decomposed into three parts:
\begin{enumerate}
  \item[(i)] $|\hat{F}_{\Ts}^{(\sstar),-1}(1-\alpha+\hat{\beta}_{\text{line}}^{(\sstar)})
    - \hat{F}_{\Ts}^{(\sstar),-1}(1-\alpha+\hat{\beta}^{(\sstar)})|$,
    where $\hat{\beta}_{\text{line}}^{(\sstar)}$ denotes the
    minimiser of~\eqref{eq:betahat_strat} over the $m$-point
    grid used in the line search (and $\hat{\beta}^{(\sstar)}$
    the continuous minimiser),
    bounded via the Lipschitz constant $K'_{\sstar}$
    and the grid spacing $\alpha/m$;
  \item[(ii)] $|\hat{F}_{\Ts}^{(\sstar),-1}(\cdot) - F_{\sstar}^{-1}(\cdot)|$
    bounded via Theorem~\ref{thm:strat_coverage};
  \item[(iii)] $|F_{\sstar}^{-1}(1-\alpha+\hat{\beta}^{(\sstar)})
    - F_{\sstar}^{-1}(1-\alpha+\beta^*)|$
    bounded via the convergence $\hat{\beta}^{(\sstar)} \to \beta^*$
    (which follows from part (ii) and the Lipschitz continuity
    of $F_{\sstar}^{-1}$).
\end{enumerate}
Adding the pointwise estimation error $|\hat{f}_{-(T+1)}(x_{T+1})
- f(x_{T+1})| \leq \deltaTs$ gives the overall bound
\eqref{eq:width_bound}, with
$M_{\sstar} = 2K_{\sstar}/|1/K'_{\sstar} - L_{\sstar}|$
as defined in \citet[proof of Theorem~3]{xu2023conformal}.\qed
